\chapter{Stratégies sur crédits hypothécaires et structurés}\label{ch:s2:mortgage-and-structured-credit-strategies}

Un propriétaire peut rembourser son prêt hypothécaire chaque fois que les taux baissent, si bien qu'une obligation hypothécaire est vendeuse d'une option que les emprunteurs risquent de mal exercer. Un trader qui modélise
mieux les emprunteurs que le marché est rémunéré pour cela, et celui qui les modélise moins bien paie. Gabaix, Krishnamurthy et Vigneron ont montré que le risque de remboursement anticipé,
qui s'annule à l'échelle agrégée, est néanmoins pris en compte dans les prix, comme si l'investisseur marginal était un spécialiste plutôt qu'un investisseur diversifié. Sur ce
chapitre, le modèle de remboursement anticipé du marché suppose un essoufflement (burnout) faible. Un strip intérêts seuls (interest-only, IO) acheté dans l'idée que l'essoufflement est fort et
couvert par des Treasuries gagne 35.8~\% de son prix si l'hypothèse se vérifie. Si les emprunteurs ne s'essoufflent jamais, il perd 25.5~\%. La construction est
\texttt{firm.mbsrv}.

\section{Vues sur le remboursement anticipé}

\begin{definition}[Trade de remboursement anticipé]\label{def:s2:mortgage-and-structured-credit-strategies:prepaytrade}
Un \emph{trade de remboursement anticipé}\index{trade de remboursement anticipé} (prepayment trade) est une position sur des titres hypothécaires dont la valeur dépend de la vitesse à laquelle les emprunteurs remboursent (pass-throughs
d'un coupon ou d'un type de pool contre un autre, strips intérêts seuls ou strips principal seul), prise parce que le modèle de remboursement anticipé du trader diffère de celui
qui figure dans les prix du marché, le risque de taux étant couvert.
\end{definition}

Le modèle du marché dans \texttt{firm.mbsrv} est celui du livre~6 (chapitre~12)~: la rotation du parc immobilier (housing turnover) suit la rampe PSA, la courbe en S du refinancement dépend de l'incitation (le
taux des emprunteurs moins le taux hypothécaire courant) avec un CPR maximal de 45~\%, et un essoufflement (burnout) qui réduit le refinancement d'un facteur $e^{-0.25 c}$, où $c$
est l'incitation positive cumulée, exprimée en points de pourcentage-années. Le pool verse à ses investisseurs 6.5~\% payés par ses emprunteurs, moins 50 points de base, et il lui reste 29~ans
à courir, et il est valorisé sur des trajectoires de taux Hull-White autour d'une courbe plate à 4~\%, avec un spread ajusté des options (OAS) de 50 points de base. Les nouveaux prêts coûtent environ 5.8~\%,
si bien que les emprunteurs ont déjà intérêt à refinancer.

Un trader estime que ces emprunteurs ont déjà renoncé à refinancer par le passé et continueront d'y renoncer~: un essoufflement de 1.0 au lieu de 0.25. Si c'est le cas, le pool
se rembourse plus lentement que ne l'anticipe le marché, et le strip intérêts seuls, qui n'est versé que tant que le capital restant dû, vaut davantage.

\section{Trades intérêts seuls et principal seul}

L'IO vaut 26.05 pour 100 de nominal du pool selon le modèle du marché. Sa valeur croît avec les taux, car des taux plus élevés ralentissent le refinancement, et la couvrir
consiste à acheter des Treasuries~: 0.51 Treasuries à dix ans de nominal 100 par IO, le ratio qui annule le risque de taux de l'IO tel que le modèle du marché le mesure
(voir le \cref{lst:s2:mortgage-and-structured-credit-strategies:io}). Le gain du trade est la valeur de l'IO selon le modèle qui se révèle juste, au
même spread ajusté des options, moins le prix payé, plus la couverture.

\begin{center}
\small
\begin{tabular}{@{}lcc@{}}
\toprule
les emprunteurs se comportent comme & gain couvert, taux inchangés & taux inférieurs de 100 pb\\
\midrule
le modèle du marché (essoufflement 0.25) & 0.0~\% & 10.1~\%\\
l'hypothèse (essoufflement 1.0) & 35.8~\% & 63.6~\%\\
refinancement plus lent (CPR maximal 33~\%) & 16.4~\% & 31.1~\%\\
aucun essoufflement (0) & $-25.5~\%$ & $-36.0~\%$\\
\bottomrule
\end{tabular}
\end{center}

Le trade est un pari sur les emprunteurs et non sur les taux, mais ce sont les taux qui déterminent l'enjeu du pari
(voir la \cref{fig:s2:mortgage-and-structured-credit-strategies:io}). Quand les taux baissent, l'incitation croît, et savoir si les emprunteurs y donnent suite est toute la
question. L'hypothèse gagne davantage, et la mauvaise hypothèse sur l'essoufflement perd davantage. Quand les taux montent, l'incitation disparaît, et tous les modèles prédisent la même
rotation lente. Même dans le modèle du marché, l'IO couvert gagne lors des grands mouvements dans un sens comme dans l'autre, parce que la couverture est dimensionnée pour les petits.

\begin{omfigure}
\begin{tikzpicture}
\begin{axis}[width=11cm,height=5.6cm,xlabel={\small déplacement parallèle de la courbe (pb)},ylabel={\small gain couvert (\% du prix de l'IO)},
  tick label style={font=\footnotesize},xmin=-150,xmax=150,ymin=-45,ymax=90,grid=major,grid style={gray!25},table/col sep=comma,
  legend style={font=\scriptsize,at={(0.5,-0.3)},anchor=north},legend columns=2]
\addplot[black,thick] table[x=shift,y=market]{figdata/strategies-2/20-mortgage-and-structured-credit-strategies/io.csv};
\addplot[omThm,thick] table[x=shift,y=view]{figdata/strategies-2/20-mortgage-and-structured-credit-strategies/io.csv};
\addplot[gray,thick,dashed] table[x=shift,y=slow]{figdata/strategies-2/20-mortgage-and-structured-credit-strategies/io.csv};
\addplot[omDef,thick] table[x=shift,y=noburnout]{figdata/strategies-2/20-mortgage-and-structured-credit-strategies/io.csv};
\legend{modèle du marché juste,hypothèse juste (essoufflement 1.0),refinancement plus lent,aucun essoufflement}
\end{axis}
\end{tikzpicture}

\omcaption{Gain d'un strip intérêts seuls acheté à la valeur du modèle du marché et couvert par des Treasuries à dix ans, en proportion de son prix, face à un
déplacement parallèle de la courbe, pour quatre comportements possibles des emprunteurs. Données~: \texttt{s2\_mbsrv.io\_scenarios}.}
\label{fig:s2:mortgage-and-structured-credit-strategies:io}
\end{omfigure}

Le strip principal seul (principal-only, PO) en est le miroir~: il gagne lorsque les remboursements sont rapides, si bien que le trader qui croit que les emprunteurs refinanceront davantage
que ne l'anticipe le marché, achète le PO. La paire de strips permet à un trader de prendre l'un ou l'autre côté d'une hypothèse sur le remboursement anticipé avec le même pool.

\subsection*{La pile des coupons}

La même hypothèse réévalue chaque coupon. Les pools dont les emprunteurs paient davantage (coupons plus élevés) ont davantage intérêt à refinancer, si bien qu'un modèle à essoufflement plus
fort les valorise davantage. Aux prix du marché, l'OAS de l'hypothèse passe de 59 points de base sur le coupon de 4.0~\% à 148 sur le coupon de 7.5~\%
(voir la \cref{fig:s2:mortgage-and-structured-credit-strategies:oas}). Avec un modèle sans essoufflement, les mêmes prix donnent des OAS qui tombent de 42 à $-326$
points de base. Un trade de valeur relative sur OAS achète le coupon au plus fort OAS selon le modèle du trader et vend celui au plus faible, couvert en duration. Le
trade gagne l'écart d'OAS seulement si le modèle est juste.

\begin{omfigure}
\begin{tikzpicture}
\begin{axis}[width=11cm,height=5.4cm,xlabel={\small coupon du pass-through (\%)},ylabel={\small OAS au prix du marché (pb)},
  tick label style={font=\footnotesize},xmin=4,xmax=7.5,ymin=-350,ymax=180,ytick={-300,-200,-100,0,100},grid=major,grid style={gray!25},table/col sep=comma,
  legend style={font=\scriptsize,at={(0.03,0.05)},anchor=south west}]
\addplot[black,thick] table[x=coupon,y=market]{figdata/strategies-2/20-mortgage-and-structured-credit-strategies/oas.csv};
\addplot[omThm,thick,mark=*,mark size=1.2pt] table[x=coupon,y=view]{figdata/strategies-2/20-mortgage-and-structured-credit-strategies/oas.csv};
\addplot[omDef,thick,mark=square*,mark size=1.2pt] table[x=coupon,y=noburnout]{figdata/strategies-2/20-mortgage-and-structured-credit-strategies/oas.csv};
\legend{modèle du marché,hypothèse (essoufflement 1.0),aucun essoufflement}
\end{axis}
\end{tikzpicture}

\omcaption{L'OAS de chaque coupon au prix que lui donne le modèle du marché (50 points de base), recalculé selon le modèle de l'hypothèse et
selon un modèle sans essoufflement. Données~: \texttt{s2\_mbsrv.stack}.}\label{fig:s2:mortgage-and-structured-credit-strategies:oas}
\end{omfigure}

Les OAS réels ne sont pas non plus plats d'un coupon à l'autre. Boyarchenko, Fuster et Lucca ont montré que les spreads des MBS présentent un net «~smile~» selon les coupons, et
l'ont attribué au risque de remboursement anticipé qui n'est pas un risque de taux, identifié à partir des prix des MBS démembrés. La variation temporelle des
spreads, en revanche, était surtout expliquée par des risques autres que le remboursement anticipé. Un smile selon les coupons est ce qu'un marché qui valorise le comportement des emprunteurs comme
un risque, et non une fonction connue, produirait~; c'est aussi ce qu'un trader disposant d'un mauvais modèle prendrait pour une erreur de valorisation.

\section{Valeur relative entre tranches}

\begin{definition}[Valeur relative entre tranches]\label{def:s2:mortgage-and-structured-credit-strategies:tranchrv}
La \emph{valeur relative entre tranches}\index{valeur relative entre tranches} (tranche relative value) consiste à échanger une tranche d'une structure de crédit mutualisée contre une autre, ou contre
l'indice ou les noms sous-jacents, parce que le modèle du trader de corrélation de défaut ou de pertes du pool diffère de celui
qui figure dans les prix des tranches.
\end{definition}

La copule gaussienne à grand pool du livre~2 (chapitre~24) valorise les tranches à partir de la probabilité de défaut du pool et d'une seule corrélation. Avec une
probabilité de défaut à cinq ans de 6~\% et un taux de recouvrement de 40~\%, la perte attendue du pool vaut 3.6~\% quelle que soit la corrélation. Les tranches se partagent cette perte de manière très différente à mesure que
la corrélation varie~:

\begin{center}
\small
\begin{tabular}{@{}lccc@{}}
\toprule
perte attendue, \% de la tranche & corrélation 0.15 & 0.30 & 0.45\\
\midrule
equity 0--3~\% & 74.45 & 59.86 & 48.11\\
mezzanine 3--7~\% & 25.32 & 24.34 & 22.11\\
mezzanine 7--10~\% & 7.52 & 11.69 & 13.00\\
senior 10--15~\% & 2.10 & 5.70 & 7.92\\
super senior 15--30~\% & 0.15 & 1.23 & 2.77\\
\bottomrule
\end{tabular}
\end{center}

La perte attendue de la tranche senior est dix-huit fois plus élevée à une corrélation de 0.45 qu'à 0.15~; celle du pool, elle, est inchangée.
Coval, Jurek et Stafford ont identifié cette fragilité comme l'une des deux caractéristiques expliquant l'essor et le déclin de la finance structurée~: des notations qui changent radicalement pour de faibles erreurs
dans l'évaluation des risques sous-jacents, et une exposition au risque systématique. Dans un article complémentaire, ils ont constaté que de nombreux produits structurés se comportaient comme des
obligations catastrophe économiques, qui ne font défaut que lors de fortes récessions, tout en offrant une rémunération bien inférieure à celle d'alternatives aux profils de gains comparables.

Un trader qui croit la corrélation égale à 0.15 alors que le marché la price à 0.30 vend de la protection sur la tranche 7--10~\% et est rémunéré à hauteur de la
perte attendue du marché, soit 11.7~\% de la tranche sur cinq ans
(le \cref{lst:s2:mortgage-and-structured-credit-strategies:tranche}). Sur 4\,000 pools simulés de 125 noms~:

\begin{center}
\small
\begin{tabular}{@{}lccc@{}}
\toprule
corrélation réelle & gain moyen (\% de la tranche) & probabilité de perte & tranche anéantie\\
\midrule
0.15 (la vue) & 3.1 & 11.3~\% & 5.4~\%\\
0.30 (le marché) & $-0.4$ & 14.5~\% & 9.4~\%\\
0.45 & $-1.4$ & 15.0~\% & 10.9~\%\\
\bottomrule
\end{tabular}
\end{center}

Même lorsque la vue est juste, le vendeur perd la tranche entière dans 5.4~\% des pools, soit environ un sur vingt. La petite perte à la corrélation propre du marché provient du caractère fini
du pool de 125 noms, que la formule à grand pool ignore. Une vue sur une tranche est une vue sur la queue de distribution, et quelques années d'historique ne suffisent pas à la confirmer.

\section{Fiches de stratégie}

\begin{strategyfile}{IO long sur une vue de remboursement anticipé lent}\label{strat:s2:mortgage-and-structured-credit-strategies:io}
\sfield{Qui paie, et pourquoi} Les détenteurs du risque de remboursement anticipé qui valorisent les emprunteurs avec un modèle générique~; les spécialistes sont les investisseurs marginaux.
\sfield{Instruments et plateformes} Strips intérêts seuls sur pools d'agences~; bons du Trésor américain ou swaps pour couvrir.
\sfield{Signal} Un modèle de remboursement anticipé fondé sur les caractéristiques du pool (taille du prêt, ancienneté, incitation passée), confronté aux vitesses du marché.
\sfield{Dimensionnement et exécution} Couvrir le risque de taux que mesure le modèle du marché~; dimensionner pour le cas où le modèle se trompe.
\sfield{Coûts} Fourchette large sur les strips~; rééquilibrage de la couverture.
\sfield{Comment elle échoue} Une vague de refinancement que le modèle sous-estime~; un changement de politique qui facilite le refinancement.
\sfield{Horizon, capacité, infrastructure} Plusieurs années~; données au niveau du prêt et modèle de remboursement anticipé.
\sfield{Backtester honnêtement} Vitesses hors échantillon par cohorte~; modèle réestimé uniquement sur les données disponibles à l'époque.
\sfield{Sources} Gabaix, Krishnamurthy et Vigneron (2007)~; ce chapitre~: $+35.8~\%$ si la vue est juste, $-25.5~\%$ sans essoufflement (burnout).
\end{strategyfile}

\begin{strategyfile}{Valeur relative OAS entre coupons}\label{strat:s2:mortgage-and-structured-credit-strategies:oas}
\sfield{Qui paie, et pourquoi} Investisseurs qui détiennent des coupons pour d'autres raisons que leur OAS~: les banques, la banque centrale, les fonds indiciels.
\sfield{Instruments et plateformes} TBA et specified pools sur l'ensemble des coupons~; dollar rolls.
\sfield{Signal} L'OAS de chaque coupon selon le modèle du trader.
\sfield{Dimensionnement et exécution} Acheter les coupons à OAS élevé, vendre ceux à OAS faible, couverts en duration.
\sfield{Coûts} Financement du roll~; payups sur les specified pools.
\sfield{Comment elle échoue} L'erreur du modèle, c'est-à-dire l'OAS~; l'offre issue des nouvelles émissions à un coupon.
\sfield{Horizon, capacité, infrastructure} Quelques mois.
\sfield{Backtester honnêtement} OAS issu du modèle tel qu'il était, sans réestimation.
\sfield{Sources} Boyarchenko, Fuster et Lucca (2019)~; ce chapitre~: un écart d'OAS de 88 points de base selon la vue.
\end{strategyfile}

\begin{strategyfile}{Valeur relative entre tranches mezzanine}\label{strat:s2:mortgage-and-structured-credit-strategies:mezz}
\sfield{Qui paie, et pourquoi} Acheteurs ou vendeurs de protection sur un niveau de la structure pour d'autres raisons que son prix~: contraintes de notation,
couverture, capital réglementaire.
\sfield{Instruments et plateformes} Tranches d'indice~; l'indice~; tranches de CLO ou de CMBS sous forme cash.
\sfield{Signal} Vues sur la corrélation ou sur les pertes du pool, confrontées aux corrélations implicites des tranches.
\sfield{Dimensionnement et exécution} Couvert en delta avec l'indice~; dimensionner pour la perte totale.
\sfield{Coûts} Fourchette large~; risque de modèle dans le ratio de couverture.
\sfield{Comment elle échoue} Une vague de défauts systémiques~; une corrélation qui saute.
\sfield{Horizon, capacité, infrastructure} Plusieurs années~; une copule et un modèle de pertes du pool.
\sfield{Backtester honnêtement} Stresser la queue de distribution~; l'historique est trop court pour confirmer une vue sur la corrélation.
\sfield{Sources} Coval, Jurek et Stafford (2009)~; ce chapitre~: gain moyen de $+3.1~\%$, tranche anéantie dans 5.4~\% des pools.
\end{strategyfile}

\section{Tutoriel~: faire mieux que le modèle du marché}\label{tut:s2:mortgage-and-structured-credit-strategies:tut}

\begin{tutorial}
\textbf{Objectif.} Valoriser un pool, ses strips et ses coupons avec un modèle de remboursement anticipé, les réévaluer avec un autre, et vendre de la protection sur une tranche à la
mauvaise corrélation. \textbf{État final~:} les trois tableaux et les deux figures.
\begin{tutsteps}
\item \textbf{Le trade IO et l'empilement des coupons}.
\omcode{code/firm/mbsrv/firm_mbsrv.py}{83}{112}{Achat de l'IO sur une vue, couvert~; OAS de chaque coupon selon la vue.}\label{lst:s2:mortgage-and-structured-credit-strategies:io}
\item \textbf{La tranche}.
\omcode{code/firm/mbsrv/firm_mbsrv.py}{115}{126}{Vente de protection à la corrélation du marché, rémunérée sur des pools à la corrélation réelle.}\label{lst:s2:mortgage-and-structured-credit-strategies:tranche}
\item \textbf{Les quatre vérités}.
\omcode{code/strategies-2/20-mortgage-and-structured-credit-strategies/python/s2_mbsrv.py}{32}{46}{Le modèle du marché, la vue et deux alternatives.}\label{lst:s2:mortgage-and-structured-credit-strategies:truths}
\item \textbf{Exécuter} \texttt{io\_summary()}, \texttt{stack()}, \texttt{tranches()}, \texttt{mezz()} et \texttt{fig\_mbsrv.py}.
\end{tutsteps}
\textbf{Pour aller plus loin.} Tirer au hasard l'essoufflement réel et mesurer la distribution du trade~; couvrir l'IO avec le modèle du marché, puis
avec celui de la vue~; couvrir en delta la tranche avec l'indice.
\end{tutorial}

\section{Construction~: valeur relative des titres hypothécaires et du crédit structuré}\label{bld:s2:mortgage-and-structured-credit-strategies:mbsrv}

\begin{build}
\bfield{Rôle} \omterm{def:s2:mortgage-and-structured-credit-strategies:prepaytrade}{Trades de remboursement anticipé} fondés sur une vue, s'appuyant sur le modèle OAS du livre~6 et le modèle de tranches du livre~2.
\bfield{Interface} \texttt{MbsConfig(\dots)}, \texttt{FlatCurve(rate)}, \texttt{strip\_paths(cfg, model, wac, oas, shift)},
\texttt{io\_trade(cfg, market, truth, side, shifts)}, \texttt{coupon\_stack(cfg, market, view, wacs)},
\texttt{tranche\_rv(p, rho\_market, rho\_true, \dots)}.
\bfield{Règles} Prix issus du modèle du marché à un OAS donné~; valeurs sous le modèle réel au même OAS~; couvertures dimensionnées avec le modèle du marché.
\bfield{Tests d'acceptation} \texttt{code/firm/mbsrv/tests/}~: la somme IO et PO égale le pass-through~; aucun gain lorsque le modèle vrai est celui du marché~; le même
modèle restitue l'OAS du marché~; des corrélations égales donnent des pertes attendues égales.
\bfield{Pour aller plus loin} Pools au niveau du prêt~; dollar rolls~; cascades de CLO.
\end{build}

\begin{omsources}
\item X.~Gabaix, A.~Krishnamurthy et O.~Vigneron, ``Limits of arbitrage: theory and evidence from the mortgage-backed securities market'',
\emph{Journal of Finance} 62(2), 2007.
\item N.~Boyarchenko, A.~Fuster et D.~O.~Lucca, ``Understanding mortgage spreads'', \emph{Review of Financial Studies} 32(10), 2019.
\item J.~Coval, J.~Jurek et E.~Stafford, ``The economics of structured finance'', \emph{Journal of Economic Perspectives} 23(1), 2009.
\item J.~Coval, J.~Jurek et E.~Stafford, ``Economic catastrophe bonds'', \emph{American Economic Review} 99(3), 2009.
\end{omsources}

\section{Exercices}

\begin{exercise}[$\star$]\label{exo:s2:mortgage-and-structured-credit-strategies:1}
Un IO acheté à 26.05 pour 100 de nominal gagne 35.8~\% de son prix. Quel est le gain pour 100 de nominal~?
\end{exercise}

\begin{exercise}[$\star$]\label{exo:s2:mortgage-and-structured-credit-strategies:2}
Sous la vue, l'OAS du coupon~7.5~\% est de 148 points de base et celui du coupon~4.0~\% de 59~; sans essoufflement, ils valent $-326$ et 42. Quelle différence d'OAS
rapporte un trade acheteur sur le 7.5~\%, vendeur sur le 4.0~\%, dans chacun des cas~?
\end{exercise}

\begin{exercise}[$\star$]\label{exo:s2:mortgage-and-structured-credit-strategies:3}
La tranche 7--10~\% d'un indice de notionnel 1 milliard a une perte attendue de 11.69~\% de la tranche. Quelle est la perte attendue en argent~?
\end{exercise}

\begin{exercise}[$\star\star$]\label{exo:s2:mortgage-and-structured-credit-strategies:4}
Pourquoi un IO perd-il de la valeur lorsque les taux baissent, et pourquoi le couvrir revient-il à acheter des bons du Trésor~?
\end{exercise}

\begin{exercise}[$\star\star$]\label{exo:s2:mortgage-and-structured-credit-strategies:5}
Pourquoi le risque de remboursement anticipé, neutre dans l'agrégat, aurait-il un prix~?
\end{exercise}

\begin{exercise}[$\star\star$]\label{exo:s2:mortgage-and-structured-credit-strategies:6}
Pourquoi la notion d'une tranche senior est-elle fragile lorsque la perte attendue du pool est bien estimée~?
\end{exercise}

\begin{exercise}[$\star\star\star$]\label{exo:s2:mortgage-and-structured-credit-strategies:7}
\emph{Programmation.} Relancer le trade sur IO avec \texttt{MbsConfig(sigma=0.012)}. Comment évoluent le prix, le gain de la vue et la perte sans essoufflement~?
\end{exercise}

\begin{exercise}[$\star\star\star$]\label{exo:s2:mortgage-and-structured-credit-strategies:8}
\emph{Trouver la faille.} «~L'OAS du coupon~7.5~\% est de 148 points de base contre 59 pour le 4.0~\%~: il est bon marché de 88 points de base, et nous le figeons.~»
\end{exercise}

\section{Problème~: mieux que le modèle du marché}

\begin{problem}[{Problème du week-end --- titres hypothécaires et crédit structuré}]\label{pb:s2:mortgage-and-structured-credit-strategies:1}
Le pool synthétique et les tranches du chapitre, ainsi que les données publiques.

\medskip\noindent\textbf{Partie I --- Le pool.}
\begin{enumerate}
\item Définir un \omterm{def:s2:mortgage-and-structured-credit-strategies:prepaytrade}{trade de remboursement anticipé}.
\item Décrire le modèle de remboursement anticipé du marché et le pool.
\item Qu'est-ce que l'essoufflement, et que suppose la vue~?
\item Qu'ont établi Gabaix, Krishnamurthy et Vigneron~?
\end{enumerate}

\medskip\noindent\textbf{Partie II --- Strips et coupons.}
\begin{enumerate}[resume]
\item Que vaut l'IO, et comment le couvre-t-on~?
\item Donner le gain couvert sous chacune des quatre hypothèses vraies.
\item Pourquoi les taux déterminent-ils l'importance de la vue~?
\item Comment la vue modifie-t-elle l'OAS de la pile de coupons~?
\end{enumerate}

\medskip\noindent\textbf{Partie III --- Tranches.}
\begin{enumerate}[resume]
\item Définir la \omterm{def:s2:mortgage-and-structured-credit-strategies:tranchrv}{valeur relative entre tranches}.
\item Comment les pertes attendues des tranches varient-elles avec la corrélation~?
\item Qu'ont établi Coval, Jurek et Stafford~?
\item Donner les résultats de la vente de la tranche 7--10~\%.
\end{enumerate}

\medskip\noindent\textbf{Partie IV --- Le verdict.}
\begin{enumerate}[resume]
\item Énoncer le \emph{résultat nommé}~: le rendement du trade sur IO lorsque la vue est juste et lorsque l'hypothèse d'essoufflement est fausse.
\item Que suggère le smile de Boyarchenko, Fuster et Lucca sur le modèle du marché~?
\item Pourquoi le gain couvert de l'IO progresse-t-il lors des grands mouvements, même sous le modèle du marché~?
\item Comment tester honnêtement un modèle de remboursement anticipé~?
\item Pour quoi faudrait-il dimensionner le trade sur IO~?
\item Pourquoi l'histoire ne peut-elle pas confirmer une vue sur la corrélation des tranches~?
\item Quelle fiche de stratégie est le plus proche des \omterm{def:s2:systematic-credit-and-bond-etf-arbitrage:factor}{facteurs crédit} du chapitre~19~?
\item En une phrase~: que vend un trader en titres hypothécaires~?
\end{enumerate}
\end{problem}

\section{Questions d'entretien}

\begin{interviewq}[$\star$ \iqroles{trader}]\label{iq:s2:mortgage-and-structured-credit-strategies:1}
Qu'est-ce que la convexité négative d'un titre adossé à des créances hypothécaires~?
\end{interviewq}

\begin{interviewq}[$\star\star$ \iqroles{chercheur}]\label{iq:s2:mortgage-and-structured-credit-strategies:2}
Comment construiriez-vous un modèle de remboursement anticipé, et quelles variables utiliseriez-vous~?
\end{interviewq}

\begin{interviewq}[$\star\star$ \iqroles{trader}]\label{iq:s2:mortgage-and-structured-credit-strategies:3}
Vous êtes acheteur d'un IO. Les taux baissent de 100 points de base. Que se passe-t-il, et que faites-vous~?
\end{interviewq}

\begin{interviewq}[$\star\star$ \iqroles{risque}]\label{iq:s2:mortgage-and-structured-credit-strategies:4}
Comment mesureriez-vous le risque d'un portefeuille de tranches au-delà de ses deltas~?
\end{interviewq}

\begin{interviewq}[$\star\star$ \iqroles{développeur}]\label{iq:s2:mortgage-and-structured-credit-strategies:5}
Le calcul d'un OAS prend une minute par pool. Comment le rendriez-vous assez rapide pour un portefeuille de dix mille pools~?
\end{interviewq}

\begin{interviewq}[$\star\star\star$ \iqroles{chercheur}]\label{iq:s2:mortgage-and-structured-credit-strategies:6}
Dans la copule gaussienne à grand pool, montrer que la perte attendue du pool ne dépend pas de la corrélation, et expliquer pourquoi celle d'une tranche en dépend.
\end{interviewq}
